\documentclass[sigconf,nonacm]{acmart}

\settopmatter{printacmref=false}
\renewcommand\footnotetextcopyrightpermission[1]{}
\pagestyle{plain}

\usepackage{booktabs}
\usepackage{amsmath}
\usepackage{array}
\usepackage{listings}
\usepackage{xcolor}

\lstset{
  basicstyle=\ttfamily\small,
  breaklines=true,
  columns=fullflexible,
  frame=single
}

\title{Multi-Initial-State Simulation in SliQSim}

\author{Daniel}
\affiliation{
  \institution{Final Project Report}
  \city{Taipei}
  \country{Taiwan}
}

\begin{document}
\sloppy

\begin{abstract}
This project extends SliQSim with simultaneous simulation of multiple labeled initial states. The method introduces Boolean selector variables that encode the initial-state label. The quantum circuit is then applied once to the combined Boolean decision diagram (BDD) representation. After the circuit is evaluated, each labeled output state is recovered by cofactoring the final BDDs with the corresponding selector assignment. The implementation supports computational-basis initial states, dense arbitrary-amplitude initial states, sparse arbitrary-amplitude initial states, repeated access to selected output labels, sequential verification, probability-tolerant comparison, BDD dumping, and a comparison script against MQT DDSIM. Experiments show that the all-in-one method produces the same support and matching probabilities as sequential SliQSim simulation on the tested circuits. Against DDSIM, the small statevector tests agree within the fixed-point precision used by SliQSim, and finite-shot sampling tests agree within expected sampling variation. For a 100-qubit sparse arbitrary-amplitude input, DDSIM cannot construct the required dense initial state vector, while the SliQSim sparse BDD input remains executable.
\end{abstract}

\maketitle

\section{Objective}
The objective is to add multi-initial-state simulation to SliQSim. Given a quantum circuit $U$ and a set of labeled initial states
\[
  \mathcal{I}=\{(\lambda_0,\lvert\psi_0\rangle),(\lambda_1,\lvert\psi_1\rangle),\ldots,(\lambda_{m-1},\lvert\psi_{m-1}\rangle)\},
\]
the simulator should evaluate all states through the same circuit without rerunning the circuit $m$ times. After this shared simulation, the user should be able to select a label $\lambda_i$ and use the existing SliQSim output functions, including sampling, statevector extraction, and probability or statistical queries.

\section{Original SliQSim State Representation}
Let the quantum register have $n$ qubits. A computational basis state is indexed by a Boolean vector
\[
  q=(q_0,q_1,\ldots,q_{n-1})\in\{0,1\}^n.
\]
SliQSim represents a quantum state as a set of Boolean functions rather than as an explicit length-$2^n$ array. For each basis assignment $q$, the amplitude is represented by SliQSim's fixed-point integer arithmetic. Internally, SliQSim stores four amplitude components and $r$ integer bits per component:
\[
  B_{c,b}(q)\in\{0,1\},
  \qquad
  c\in\{0,1,2,3\},
  \qquad
  b\in\{0,1,\ldots,r-1\}.
\]
The implementation stores these roots in
\[
  \texttt{All\_Bdd[c][b]}.
\]
Each root is a BDD over the qubit variables. The integer value of one component is reconstructed from its bit slices. SliQSim also maintains a global scaling parameter $k$, which records powers of $\sqrt{2}$ introduced by gates such as Hadamard and rotations. Therefore the simulator does not store one floating-point complex number per basis state. It stores bit-sliced Boolean functions whose evaluations reconstruct the amplitude under SliQSim's original arithmetic model.

This project does not replace the original gate arithmetic. All gate functions still update the same \texttt{All\_Bdd[c][b]} roots by BDD operations such as cofactor, conjunction, disjunction, and exclusive-or.

\section{Multi-Initial-State Encoding}
\subsection{Selector Variables}
For $m$ initial states, the implementation introduces
\[
  \ell=\lceil \log_2 m\rceil
\]
selector variables
\[
  s=(s_0,s_1,\ldots,s_{\ell-1})\in\{0,1\}^{\ell}.
\]
These variables are allocated after the qubit variables in the CUDD manager. If the circuit has $n$ qubits, then qubit variables use indices $0,\ldots,n-1$, and selector variables use indices
\[
  n,\ldots,n+\ell-1.
\]
The CUDD manager is initialized with
\[
  n+\ell
\]
primary BDD variables. For label index $i$, let $\operatorname{bin}_{\ell}(i)$ be the $\ell$-bit binary representation of $i$. Define the selector cube
\[
  C_i(s)=
  \bigwedge_{j=0}^{\ell-1}
  \begin{cases}
    s_j, & \operatorname{bin}_{\ell}(i)_j=1,\\
    \neg s_j, & \operatorname{bin}_{\ell}(i)_j=0.
  \end{cases}
\]
This cube is true only for the selector assignment corresponding to label $\lambda_i$.

\subsection{Combined Initial State}
Let $\alpha_i(q)$ be the amplitude of basis state $q$ in initial state $\lvert\psi_i\rangle$. The combined multi-initial-state amplitude function is
\[
  \Psi(q,s)=
  \begin{cases}
    \alpha_i(q), & s=\operatorname{bin}_{\ell}(i),\\
    0, & s \notin \{\operatorname{bin}_{\ell}(0),\ldots,\operatorname{bin}_{\ell}(m-1)\}.
  \end{cases}
\]
Equivalently,
\[
  \Psi(q,s)=
  \sum_{i=0}^{m-1} C_i(s)\alpha_i(q).
\]
For each fixed-point component $c$ and bit $b$, the BDD root is constructed as the union of all labeled terms whose encoded bit is one:
\[
  B^{\mathrm{multi}}_{c,b}(q,s)
  =
  \bigvee_{i=0}^{m-1}
  \left(C_i(s)\wedge B^{(i)}_{c,b}(q)\right).
\]
This is the central construction. It stores all labels in one set of BDD roots, but the labels remain disjoint because the selector cubes $C_i$ are mutually exclusive.

\subsection{Basis Initial States}
For a basis state string $x_i\in\{0,1\}^n$, the amplitude is one on that assignment and zero elsewhere:
\[
  \alpha_i(q)=
  \begin{cases}
    1, & q=x_i,\\
    0, & q\ne x_i.
  \end{cases}
\]
The BDD term inserted into the real component is
\[
  C_i(s)\wedge
  \bigwedge_{j=0}^{n-1}
  \begin{cases}
    q_j, & (x_i)_j=1,\\
    \neg q_j, & (x_i)_j=0.
  \end{cases}
\]
The value one is encoded using SliQSim's fixed-point integer representation.

\subsection{Dense Arbitrary-Amplitude Initial States}
A dense arbitrary-amplitude line gives all $2^n$ amplitudes in computational-basis order:
\[
  \lambda_i\quad a_0\quad a_1\quad \ldots\quad a_{2^n-1}.
\]
Each value $a_t\in\mathbb{C}$ is parsed as
\[
  a_t=\operatorname{Re}(a_t)+i\operatorname{Im}(a_t).
\]
The implementation encodes real and imaginary parts by fixed-point scaling. Let the number of fractional bits be
\[
  f=\min(24,r-2).
\]
Then a real number $x$ is encoded as
\[
  \widehat{x}=\operatorname{round}(x2^f).
\]
The simulator sets
\[
  k=2f,
\]
so that the stored integer is interpreted with scale $2^{-f}$. Therefore the decoded value is approximately
\[
  x\approx \widehat{x}2^{-f}.
\]
This approximation is the source of small probability differences against floating-point DDSIM statevectors.

\subsection{Sparse Arbitrary-Amplitude Initial States}
Dense input is impossible for large $n$. For example, $n=100$ would require $2^{100}$ amplitudes. Therefore the implementation also supports sparse entries:
\[
  \lambda_i\quad x_{i,0}:a_{i,0}\quad x_{i,1}:a_{i,1}\quad \ldots
\]
where each $x_{i,t}\in\{0,1\}^n$ is a basis string and $a_{i,t}\in\mathbb{C}$ is its amplitude. The BDD contribution for each sparse term is
\[
  C_i(s)\wedge Q_{i,t}(q),
\]
where $Q_{i,t}$ is the basis cube for $x_{i,t}$. Only nonzero terms are inserted. This is the input format used for the 100-qubit Bernstein--Vazirani experiment.

\section{Why Gate Implementations Do Not Change}
The selector variables are not quantum variables. They are label variables. Every quantum gate in SliQSim acts on qubit variables only. A gate $G$ on the quantum register is applied to the combined state as
\[
  I_{\mathrm{selector}}\otimes G.
\]
Thus, if the initial combined state is
\[
  \Psi(q,s)=\sum_i C_i(s)\alpha_i(q),
\]
then after applying circuit $U$ the result is
\[
  \Psi'(q,s)
  =
  \sum_i C_i(s)(U\alpha_i)(q).
\]
The selector cube $C_i(s)$ is unchanged.

At the BDD level, the original gate code performs operations such as
\[
  \operatorname{Cofactor}(B_{c,b},q_t),
  \quad
  \operatorname{Cofactor}(B_{c,b},\neg q_t),
  \quad
  B_{c,b}\wedge q_t,
  \quad
  B_{c,b}\vee q_t,
  \quad
  B_{c,b}\oplus q_t.
\]
The variable $q_t$ is always a qubit variable with index less than $n$. The selector variables have indices at least $n$. Therefore the gate code transforms the dependence on $q$ while preserving the dependence on $s$. This is why no gate-specific rewrite is needed.

\section{Recovering One Output State}
After the circuit has been simulated once, the user selects a label $\lambda_i$. The implementation constructs the selector cube $C_i(s)$ and cofactors every BDD root:
\[
  B^{(i),\mathrm{out}}_{c,b}(q)
  =
  \operatorname{Cofactor}\left(B^{\mathrm{multi,out}}_{c,b}(q,s), C_i(s)\right).
\]
Because the selector cubes are disjoint, this operation extracts exactly the output state for label $\lambda_i$:
\[
  \Psi'_i(q)=(U\alpha_i)(q).
\]
After this projection, SliQSim's existing functions can be called without modification. The implementation temporarily replaces \texttt{All\_Bdd} with the selected cofactored roots, performs sampling, statevector extraction, or query evaluation, and then restores the original multi-state roots. This permits repeated access to different labels without rerunning the quantum circuit.

\section{Merged BDD for Measurement}
The original SliQSim measurement procedure merges the bit-sliced roots into one large BDD called \texttt{bigBDD}. With $r$ integer bits and $w=4$ amplitude components, the merged BDD needs auxiliary variables to select one component and one integer bit. The number of auxiliary variables is
\[
  \lceil\log_2 w\rceil+\lceil\log_2 r\rceil.
\]
With the default $w=4$ and $r=32$, this is
\[
  \lceil\log_2 4\rceil+\lceil\log_2 32\rceil=2+5=7.
\]
These auxiliary variables are not qubits and not initial-state selectors. They only encode which amplitude component and which fixed-point bit are being traversed during probability computation.

\section{Verification Method}
\subsection{Sequential SliQSim Verification}
The implementation includes a sequential verification mode. For each selected label $\lambda_i$, it constructs a simulator containing only that one initial state, runs the same circuit, and records the output. This produces the reference set
\[
  \{U\lvert\psi_i\rangle\}_{i\in S},
\]
where $S$ is the selected label set. The all-in-one output is compared against this sequential output.

For statevectors, the comparison is performed by converting amplitudes to probabilities. For sampling results, counts are first normalized:
\[
  \widehat{p}_{i,y}=\frac{N_{i,y}}{N_i},
\]
where $N_{i,y}$ is the number of times label $i$ produced output bitstring $y$, and $N_i$ is the shot count for label $i$.

\subsection{Support Check}
Before comparing numeric probabilities, the verifier checks support equality. For each label $i$, define
\[
  \operatorname{supp}(p_i)=\{y:p_{i,y}>0\}.
\]
The support check requires
\[
  \operatorname{supp}(p_i^{\mathrm{multi}})
  =
  \operatorname{supp}(p_i^{\mathrm{seq}})
\]
for all selected labels. If an output bitstring appears in one mode but not the other, the report identifies the missing label and missing output.

\subsection{Probability Difference}
After the support check, the maximum absolute probability difference is computed:
\[
  \Delta_{\max}
  =
  \max_{i,y}
  \left|
    p_{i,y}^{\mathrm{multi}}
    -
    p_{i,y}^{\mathrm{reference}}
  \right|.
\]
For finite-shot sampling, nonzero differences are expected because the two simulators draw random samples independently. Therefore the verification command accepts a tolerance. In the report experiments, statevector comparisons show differences below $10^{-6}$, while sampling comparisons show differences on the order expected from $10{,}000$ shots.

\section{Implemented User Interface}
The implementation adds the following user-facing options:
\begin{itemize}
  \item \texttt{--init\_states FILE}: read labeled initial states.
  \item \texttt{--select\_init LABELS}: access all labels or a comma-separated subset.
  \item \texttt{--sequential\_init}: run selected labels one by one.
  \item \texttt{--verify\_init}: run all-in-one and sequential modes and compare them.
  \item \texttt{--verify\_tol VALUE}: probability tolerance for verification.
  \item \texttt{--dump\_bdds PREFIX}: export Graphviz DOT files for initial BDD slices and final merged BDDs.
\end{itemize}

The comparison script is:
\[
  \texttt{tools/compare\_simulators.py}.
\]
It runs SliQSim all-in-one, SliQSim sequential, and MQT DDSIM sequential simulation. DDSIM is run label by label because it does not implement the selector-variable encoding. Each DDSIM label is executed in a child process with a hard wall-time limit \texttt{--ddsim-timeout}. If DDSIM cannot construct a dense initial state, or if the label exceeds the timeout, the label is reported as aborted.

\section{Experimental Setup}
All experiments were run on the same local workstation using the current SliQSim executable in this repository. The random seed was fixed to zero. Unless otherwise stated, SliQSim used the default integer bit size $r=32$. Sampling experiments used $10{,}000$ shots. DDSIM was called through the Python package \texttt{mqt.ddsim}. The reported peak memory values are resident-set-size measurements. SliQSim also reports its own peak memory and maximum BDD node count.

\subsection{Test Cases}
The report uses the following circuits and initial-state files:
\begin{itemize}
  \item \texttt{examples/bell\_state.qasm} with \texttt{init\_state/init\_states.txt}: two basis labels.
  \item \texttt{examples/bell\_state\_measure.qasm} with \texttt{init\_state/init\_states.txt}: two basis labels and measurement.
  \item \texttt{examples/report\_entangle3.qasm} with \texttt{examples/report\_basis4\_init.txt}: four basis labels on a 3-qubit circuit.
  \item \texttt{examples/report\_entangle3.qasm} with \texttt{examples/report\_sparse\_complex\_init.txt}: four sparse complex arbitrary labels.
  \item \texttt{examples/report\_entangle3\_measure.qasm} with \texttt{examples/report\_sparse\_complex\_init.txt}: measured sparse complex arbitrary labels.
  \item \texttt{examples/bv100.qasm} with \texttt{examples/bv100\_sparse\_arbitrary\_init.txt}: three selected sparse arbitrary labels from a 100-qubit benchmark.
\end{itemize}

\section{Experimental Results}
\subsection{Runtime and Size Statistics}
\begin{table*}[t]
\centering
\caption{Runtime, memory, and BDD size statistics. ``Multi'' is SliQSim all-in-one multi-initial-state simulation. ``Seq.'' is SliQSim sequential simulation. ``DDSIM'' is MQT DDSIM sequential simulation.}
\label{tab:runtime}
\begin{tabular}{llrrrrrr}
\toprule
Case & Mode & Wall s & RSS bytes & SliQSim s & SliQSim mem & Gates & Max nodes \\
\midrule
Bell statevector & Multi & 0.019646 & 3,305,472 & 0.008155 & 18,219,008 & 2 & 6 \\
Bell statevector & Seq. & 0.021139 & 7,757,824 & 0.013330 & 19,660,800 & -- & -- \\
Bell statevector & DDSIM & 0.182977 & 114,528,256 & -- & -- & -- & -- \\
\midrule
Bell sampling & Multi & 0.030328 & 11,702,272 & 0.025338 & 18,087,936 & 2 & 21 \\
Bell sampling & Seq. & 0.041724 & 13,799,424 & 0.031382 & 19,529,728 & -- & -- \\
Bell sampling & DDSIM & 0.172827 & 114,704,384 & -- & -- & -- & -- \\
\midrule
3q basis statevector & Multi & 0.020347 & 11,567,104 & 0.009299 & 18,309,120 & 6 & 23 \\
3q basis statevector & Seq. & 0.021183 & 6,721,536 & 0.016246 & 19,488,768 & -- & -- \\
3q basis statevector & DDSIM & 0.358148 & 114,315,264 & -- & -- & -- & -- \\
\midrule
3q complex statevector & Multi & 0.019701 & 3,330,048 & 0.007975 & 18,337,792 & 6 & 65 \\
3q complex statevector & Seq. & 0.021281 & 5,468,160 & 0.016041 & 19,779,584 & -- & -- \\
3q complex statevector & DDSIM & 0.446059 & 114,520,064 & -- & -- & -- & -- \\
\midrule
3q complex sampling & Multi & 0.081051 & 13,500,416 & 0.068653 & 18,350,080 & 6 & 138 \\
3q complex sampling & Seq. & 0.082979 & 13,582,336 & 0.075722 & 19,791,872 & -- & -- \\
3q complex sampling & DDSIM & 0.782672 & 114,597,888 & -- & -- & -- & -- \\
\midrule
BV100 sparse sampling & Multi & 4.091575 & 17,227,776 & 4.075810 & 18,317,312 & 299 & 4,272 \\
BV100 sparse sampling & Seq. & 0.626792 & 37,425,152 & 0.622212 & 37,425,152 & -- & -- \\
BV100 sparse sampling & DDSIM & 0.083774 & 114,606,080 & -- & -- & -- & -- \\
\bottomrule
\end{tabular}
\end{table*}

\subsection{Correctness and Probability Agreement}
\begin{table*}[t]
\centering
\caption{Probability comparison results. Support match means that both compared modes produced the same set of output bitstrings for every selected label.}
\label{tab:probability}
\begin{tabular}{llccrl}
\toprule
Case & Comparison & Same labels & Support match & $\Delta_{\max}$ & Max label \\
\midrule
Bell statevector & Multi vs Seq. & true & true & 0 & -- \\
Bell statevector & Multi vs DDSIM & true & true & $3.09\times10^{-7}$ & one \\
\midrule
Bell sampling & Multi vs Seq. & true & true & 0.0006 & one \\
Bell sampling & Multi vs DDSIM & true & true & 0.0028 & zero \\
\midrule
3q basis statevector & Multi vs Seq. & true & true & 0 & -- \\
3q basis statevector & Multi vs DDSIM & true & true & $5.52\times10^{-7}$ & high \\
\midrule
3q complex statevector & Multi vs Seq. & true & true & 0 & -- \\
3q complex statevector & Multi vs DDSIM & true & true & $3.68\times10^{-7}$ & four\_term \\
\midrule
3q complex sampling & Multi vs Seq. & true & true & 0.0125 & four\_term \\
3q complex sampling & Multi vs DDSIM & true & true & 0.0161 & phi\_plus \\
\midrule
BV100 sparse sampling & Multi vs Seq. & true & true & 0.0450 & arb01 \\
BV100 sparse sampling & Multi vs DDSIM & false & false & 0.2790 & arb02 \\
\bottomrule
\end{tabular}
\end{table*}

The statevector comparisons have differences below $10^{-6}$. These differences are consistent with SliQSim's fixed-point encoding of decimal amplitudes. The sampling comparisons have larger differences because the compared simulators perform finite-shot random sampling. With $10{,}000$ shots, differences on the order of $10^{-2}$ are plausible for probabilities near $1/2$ or $1/4$.

For BV100, DDSIM did not produce reference probabilities because the arbitrary sparse 100-qubit states would need to be expanded into dense vectors of length $2^{100}$. The comparison script therefore reports the DDSIM labels as aborted. The SliQSim multi-vs-sequential comparison still has matching labels and matching support. Its maximum sampled probability difference is 0.045 for 1,000 shots, which is a finite-shot sampling difference rather than a statevector disagreement.

\subsection{DDSIM Per-Label Statistics}
\begin{table*}[t]
\centering
\caption{DDSIM per-label summary for representative non-aborted cases. Backend time is the time reported by DDSIM. Total wall time also includes Python, Qiskit circuit construction, and process overhead.}
\label{tab:ddsim}
\begin{tabular}{llrrrr}
\toprule
Case & Label set & Total wall s per label & Backend s per label & Depth & Operations \\
\midrule
Bell statevector & zero, one & 0.085--0.086 & 0.041--0.047 & 2--3 & 2--3 \\
Bell sampling & zero, one & 0.085--0.086 & 0.042 & 3--4 & 4--5 \\
3q basis statevector & four labels & 0.075--0.096 & 0.043--0.053 & 5--6 & 6--9 \\
3q complex statevector & four labels & 0.105--0.116 & 0.045--0.048 & 15 & 20 \\
3q complex sampling & four labels & 0.188--0.199 & 0.130--0.143 & 16 & 23 \\
\bottomrule
\end{tabular}
\end{table*}

\section{Interpretation}
The experiments support four conclusions.

First, the selector-variable construction is functionally correct for the tested circuits. SliQSim all-in-one and SliQSim sequential simulation always had the same selected label set and the same output support. For statevector experiments, their probability distributions were identical after parsing. For sampling experiments, the differences were finite-shot sampling differences.

Second, the gate implementations do not need to be modified. Since the gate code only cofactors and combines on qubit variables, selector variables are preserved automatically. This realizes the operator $I_{\mathrm{selector}}\otimes U$.

Third, sparse arbitrary-amplitude input is necessary for large circuits. DDSIM can initialize small arbitrary states by expanding the vector and prepending an initialization circuit. For 100 qubits, dense expansion is infeasible. SliQSim can still construct the sparse BDD terms because it only inserts the nonzero basis cubes.

Fourth, all-in-one multi-initial-state simulation is not guaranteed to be faster than sequential simulation for every input. The all-in-one BDD includes selector variables and may increase BDD size. It is most useful when the selected initial states share BDD structure during the circuit. When there is little sharing or when only a few labels are selected, sequential simulation can be faster. The benefit is therefore structural, not automatic.

\section{Limitations}
The implementation has the following limitations:
\begin{itemize}
  \item Decimal arbitrary amplitudes are fixed-point approximations controlled by \texttt{--r}. They are not exact real numbers.
  \item Dense arbitrary-amplitude input is only practical for small $n$ because it requires $2^n$ entries.
  \item BDD traversal is polynomial in the number of BDD nodes, but the BDD size itself can be exponential in the number of qubits.
  \item DDSIM is not an all-in-one multi-initial-state baseline. It is used as a sequential external reference.
  \item Finite-shot sampling comparisons are random experiments. Probability differences decrease statistically as the shot count increases, but exact equality is not expected.
\end{itemize}

\section{Reproducibility Commands}
The main commands used to generate the tables are shown below.

\begin{lstlisting}[language=bash]
python3 tools/compare_simulators.py \
  --qasm examples/bell_state.qasm \
  --init-states init_state/init_states.txt \
  --type 1 --seed 0 --ddsim-timeout 10 \
  --json-out /tmp/report_bell_basis.json

python3 tools/compare_simulators.py \
  --qasm examples/bell_state_measure.qasm \
  --init-states init_state/init_states.txt \
  --type 0 --shots 10000 --seed 0 --ddsim-timeout 10 \
  --json-out /tmp/report_bell_measure.json

python3 tools/compare_simulators.py \
  --qasm examples/report_entangle3.qasm \
  --init-states examples/report_basis4_init.txt \
  --type 1 --seed 0 --ddsim-timeout 10 \
  --json-out /tmp/report_entangle3_basis.json

python3 tools/compare_simulators.py \
  --qasm examples/report_entangle3.qasm \
  --init-states examples/report_sparse_complex_init.txt \
  --type 1 --seed 0 --ddsim-timeout 10 \
  --json-out /tmp/report_entangle3_complex.json

python3 tools/compare_simulators.py \
  --qasm examples/report_entangle3_measure.qasm \
  --init-states examples/report_sparse_complex_init.txt \
  --type 0 --shots 10000 --seed 0 --ddsim-timeout 10 \
  --json-out /tmp/report_entangle3_complex_measure.json

python3 tools/compare_simulators.py \
  --qasm examples/bv100.qasm \
  --init-states examples/bv100_sparse_arbitrary_init.txt \
  --select-init arb00,arb01,arb02 \
  --type 0 --shots 1000 --seed 0 --ddsim-timeout 1 \
  --json-out /tmp/report_bv100_sparse3.json
\end{lstlisting}

\section{Conclusion}
This project implements multi-initial-state simulation in SliQSim by adding selector variables to the BDD representation of the initial state. The mathematical construction is a direct sum over mutually exclusive selector cubes. The original gate operations remain valid because each gate acts only on quantum variables, so the selector variables are preserved as labels. Output states are recovered by cofactoring with the selected label cube. The implementation also adds arbitrary-amplitude input, sparse input, label selection, sequential verification, probability-tolerant comparison, BDD dumping, and DDSIM comparison tooling. The experimental results show agreement between all-in-one and sequential SliQSim simulation, close agreement with DDSIM on small cases, and a clear scalability advantage of sparse BDD initialization over DDSIM dense initialization for the 100-qubit arbitrary-state benchmark.

\end{document}
